\mxunit{premiere-matrice}{Matrice de couverture du programme}\label{MATRICE}
Les attendus sont reformulés et décomposés en entrées vérifiables. « Couvert » indique une ressource localisée, et non une maîtrise attestée de l’élève. Les approfondissements facultatifs du BO ne sont pas tous imposés. Les ressources de seconde sont entretenues dans R001–R030 et STATS.
\small\begin{longtable}{p{.1\linewidth}p{.43\linewidth}p{.19\linewidth}p{.16\linewidth}}\toprule ID & Attendu / type & Cours, exemple & Exercice \\\midrule\endhead
BO001 & Appartenance et inclusion (Contenu) & LOGIQUE, p.\,\pageref{LOGIQUE} ; LOGIQUE & LOG-EX, p.\,\pageref{LOG-EX} \\
BO002 & Ensemble vide, réunion, intersection et complémentaire (Contenu) & LOGIQUE, p.\,\pageref{LOGIQUE} ; LOGIQUE & LOG-EX, p.\,\pageref{LOG-EX} \\
BO003 & Ensembles de nombres et intervalles (Contenu) & LOGIQUE, p.\,\pageref{LOGIQUE} ; LOGIQUE & LOG-EX, p.\,\pageref{LOG-EX} \\
BO004 & Couple et produit cartésien (Contenu) & LOGIQUE, p.\,\pageref{LOGIQUE} ; LOGIQUE & LOG-EX, p.\,\pageref{LOG-EX} \\
BO005 & Cardinal d’un ensemble fini (Contenu) & LOGIQUE, p.\,\pageref{LOGIQUE} ; LOGIQUE & LOG-EX, p.\,\pageref{LOG-EX} \\
BO006 & Connecteurs et/ou (Capacité) & LOGIQUE, p.\,\pageref{LOGIQUE} ; LOGIQUE & LOG-EX, p.\,\pageref{LOG-EX} \\
BO007 & Contre-exemple (Capacité) & LOGIQUE, p.\,\pageref{LOGIQUE} ; LOGIQUE & LOG-EX, p.\,\pageref{LOG-EX} \\
BO008 & Implication et équivalence (Capacité) & LOGIQUE, p.\,\pageref{LOGIQUE} ; LOGIQUE & LOG-EX, p.\,\pageref{LOG-EX} \\
BO009 & Réciproque et contraposée (Capacité) & LOGIQUE, p.\,\pageref{LOGIQUE} ; LOGIQUE & LOG-EX, p.\,\pageref{LOG-EX} \\
BO010 & Condition nécessaire et suffisante (Capacité) & LOGIQUE, p.\,\pageref{LOGIQUE} ; LOGIQUE & LOG-EX, p.\,\pageref{LOG-EX} \\
BO011 & Identité, équation, variable, inconnue et paramètre (Capacité) & LOGIQUE, p.\,\pageref{LOGIQUE} ; LOGIQUE & LOG-EX, p.\,\pageref{LOG-EX} \\
BO012 & Quantificateurs et quantification implicite (Capacité) & LOGIQUE, p.\,\pageref{LOGIQUE} ; LOGIQUE & LOG-EX, p.\,\pageref{LOG-EX} \\
BO013 & Négation des propositions quantifiées (Capacité) & LOGIQUE, p.\,\pageref{LOGIQUE} ; LOGIQUE & LOG-EX, p.\,\pageref{LOG-EX} \\
BO014 & Raisonnement par disjonction des cas (Capacité) & LOGIQUE, p.\,\pageref{LOGIQUE} ; LOGIQUE & LOG-EX, p.\,\pageref{LOG-EX} \\
BO015 & Raisonnement par l’absurde (Capacité) & LOGIQUE, p.\,\pageref{LOGIQUE} ; LOGIQUE & LOG-EX, p.\,\pageref{LOG-EX} \\
BO016 & Raisonnement par contraposée (Capacité) & LOGIQUE, p.\,\pageref{LOGIQUE} ; LOGIQUE & LOG-EX, p.\,\pageref{LOG-EX} \\
BO017 & Variables et affectation (Capacité) & PYTHON, p.\,\pageref{PYTHON} ; PYTHON & TP, p.\,\pageref{TP} \\
BO018 & Fonctions et retour de valeur (Capacité) & PYTHON, p.\,\pageref{PYTHON} ; PYTHON & TP, p.\,\pageref{TP} \\
BO019 & Boucles bornées et non bornées avec arrêt (Capacité) & PYTHON, p.\,\pageref{PYTHON} ; PYTHON & TP, p.\,\pageref{TP} \\
BO020 & Instructions conditionnelles (Capacité) & PYTHON, p.\,\pageref{PYTHON} ; PYTHON & TP, p.\,\pageref{TP} \\
BO021 & Liste en extension, ajouts successifs et compréhension (Capacité) & PYTHON, p.\,\pageref{PYTHON} ; PYTHON & TP, p.\,\pageref{TP} \\
BO022 & Manipulation, ajout et suppression des éléments et indices (Capacité) & PYTHON, p.\,\pageref{PYTHON} ; PYTHON & TP, p.\,\pageref{TP} \\
BO023 & Parcours d’une liste et itération sur ses éléments (Capacité) & PYTHON, p.\,\pageref{PYTHON} ; PYTHON & TP, p.\,\pageref{TP} \\
BO024 & Taux direct (Capacité) & R006, p.\,\pageref{R006} ; R006 & R006, p.\,\pageref{R006} \\
BO025 & Valeur initiale (Capacité) & R007, p.\,\pageref{R007} ; R007 & R007, p.\,\pageref{R007} \\
BO026 & Taux d’évolution en pourcentage (Capacité) & R008, p.\,\pageref{R008} ; R008 & R008, p.\,\pageref{R008} \\
BO027 & Évolutions successives (Capacité) & R009, p.\,\pageref{R009} ; R009 & R009, p.\,\pageref{R009} \\
BO028 & Évolution réciproque (Capacité) & R010, p.\,\pageref{R010} ; R010 & R010, p.\,\pageref{R010} \\
BO029 & Produit nul (Capacité) & R011, p.\,\pageref{R011} ; R011 & R011, p.\,\pageref{R011} \\
BO030 & Signe d’expression affine (Capacité) & R014, p.\,\pageref{R014} ; R014 & R014, p.\,\pageref{R014} \\
BO031 & Signe d’expression factorisée du second degré (Capacité) & C02-U4, p.\,\pageref{C02-U4} ; C02-EX4 & C02-E1, p.\,\pageref{C02-E1} \\
BO032 & Développer, factoriser et réduire (Capacité) & R012, p.\,\pageref{R012} ; R012 & R012, p.\,\pageref{R012} \\
BO033 & Équations et inéquations graphiques (Capacité) & R078, p.\,\pageref{R078} ; R078 & R078, p.\,\pageref{R078} \\
BO034 & Signe et variations graphiques (Capacité) & R078, p.\,\pageref{R078} ; R078 & R078, p.\,\pageref{R078} \\
BO035 & Tracer une droite par équation ou point et pente (Capacité) & R016, p.\,\pageref{R016} ; R016 & R016, p.\,\pageref{R016} \\
BO036 & Lire l’équation d’une droite (Capacité) & R015, p.\,\pageref{R015} ; R015 & R015, p.\,\pageref{R015} \\
BO037 & Coefficient directeur par deux points (Capacité) & R015, p.\,\pageref{R015} ; R015 & R015, p.\,\pageref{R015} \\
BO038 & Graphiques, histogrammes, barres, circulaires et boîtes (Capacité) & STATS, p.\,\pageref{STATS} ; STATS & STATS-EX, p.\,\pageref{STATS-EX} \\
BO039 & Passer des données à une représentation (Capacité) & STATS, p.\,\pageref{STATS} ; STATS & STATS-EX, p.\,\pageref{STATS-EX} \\
BO040 & Indicateurs statistiques (Capacité) & STATS, p.\,\pageref{STATS} ; STATS & STATS-EX, p.\,\pageref{STATS-EX} \\
BO041 & Conditionnelles en tableau ou arbre (Capacité) & C05-U1, p.\,\pageref{C05-U1} ; C05-EX1 & C05-E1, p.\,\pageref{C05-E1} \\
BO042 & Distinguer intersection et deux conditionnements (Capacité) & C05-U1, p.\,\pageref{C05-U1} ; C05-EX1 & C05-E1, p.\,\pageref{C05-E1} \\
BO043 & Modes explicite, récurrent, algorithmique et motifs (Contenu) & C01-U1, p.\,\pageref{C01-U1} ; C01-EX1 & C01-E1, p.\,\pageref{C01-E1} \\
BO044 & Notations de la suite et d’un terme (Contenu) & C01-U1, p.\,\pageref{C01-U1} ; C01-EX1 & C01-E1, p.\,\pageref{C01-E1} \\
BO045 & Passer entre langue naturelle, algèbre et graphique (Capacité) & C01-U1, p.\,\pageref{C01-U1} ; C01-EX1 & C01-E5, p.\,\pageref{C01-E5} \\
BO046 & Modéliser un motif ou un comptage (Capacité) & C01-U1, p.\,\pageref{C01-U1} ; C01-EX1 & C01-E5, p.\,\pageref{C01-E5} \\
BO047 & Calculer des termes explicites, récurrents ou algorithmiques (Capacité) & C01-U1, p.\,\pageref{C01-U1} ; C01-EX1 & C01-E5, p.\,\pageref{C01-E5} \\
BO048 & Suite arithmétique, terme général, accroissement constant, lien affine (Contenu) & C01-U2, p.\,\pageref{C01-U2} ; C01-EX2 & C01-E2, p.\,\pageref{C01-E2} \\
BO049 & Suite géométrique, terme général, taux constant, lien exponentiel (Contenu) & C01-U3, p.\,\pageref{C01-U3} ; C01-EX3 & C01-E3, p.\,\pageref{C01-E3} \\
BO050 & Terme général d’une suite arithmétique (Démonstration) & C01-U2, p.\,\pageref{C01-U2} ; C01-EX2 & C01-E8, p.\,\pageref{C01-E8} \\
BO051 & Terme général d’une suite géométrique (Démonstration) & C01-U3, p.\,\pageref{C01-U3} ; C01-EX3 & C01-E3, p.\,\pageref{C01-E3} \\
BO052 & Somme des premiers entiers (Démonstration) & C01-U4, p.\,\pageref{C01-U4} ; C01-EX4 & C01-E4, p.\,\pageref{C01-E4} \\
BO053 & Somme de 1 à q puissance n (Démonstration) & C01-U4, p.\,\pageref{C01-U4} ; C01-EX4 & C01-E4, p.\,\pageref{C01-E4} \\
BO054 & Calculer des sommes de termes consécutifs (Capacité) & C01-U4, p.\,\pageref{C01-U4} ; C01-EX4 & C01-E4, p.\,\pageref{C01-E4} \\
BO055 & Déterminer le sens de variation (Capacité) & C01-U3, p.\,\pageref{C01-U3} ; C01-EX3 & C01-E6, p.\,\pageref{C01-E6} \\
BO056 & Modéliser une croissance linéaire ou exponentielle discrète (Capacité) & C01-U3, p.\,\pageref{C01-U3} ; C01-EX3 & C01-E6, p.\,\pageref{C01-E6} \\
BO057 & Limite intuitive finie, infinie ou absence de limite (Contenu) & C01-U5, p.\,\pageref{C01-U5} ; C01-EX5 & C01-E6, p.\,\pageref{C01-E6} \\
BO058 & Conjecturer une limite dans des cas simples (Capacité) & C01-U5, p.\,\pageref{C01-U5} ; C01-EX5 & C01-AP5, p.\,\pageref{C01-AP5} \\
BO059 & Formes factorisée, développée et canonique (Contenu) & C02-U1, p.\,\pageref{C02-U1} ; C02-EX1 & C02-E1, p.\,\pageref{C02-E1} \\
BO060 & Racines et choix de forme (Contenu) & C02-U1, p.\,\pageref{C02-U1} ; C02-EX1 & C02-E1, p.\,\pageref{C02-E1} \\
BO061 & Somme et produit des racines (Contenu) & C02-U3, p.\,\pageref{C02-U3} ; C02-EX3 & C02-E3, p.\,\pageref{C02-E3} \\
BO062 & Discriminant, résolution et factorisation éventuelle (Contenu) & C02-U2, p.\,\pageref{C02-U2} ; C02-EX2 & C02-E3, p.\,\pageref{C02-E3} \\
BO063 & Étudier le signe d’une forme factorisée (Capacité) & C02-U4, p.\,\pageref{C02-U4} ; C02-EX4 & C02-E5, p.\,\pageref{C02-E5} \\
BO064 & Choisir une forme pour équation, inéquation, optimisation ou variations (Capacité) & C02-U4, p.\,\pageref{C02-U4} ; C02-EX4 & C02-E5, p.\,\pageref{C02-E5} \\
BO065 & Déterminer les polynômes s’annulant en deux réels distincts (Capacité) & C02-U3, p.\,\pageref{C02-U3} ; C02-EX3 & C02-E4, p.\,\pageref{C02-E4} \\
BO066 & Diversifier les stratégies de factorisation (Capacité) & C02-U3, p.\,\pageref{C02-U3} ; C02-EX3 & C02-E4, p.\,\pageref{C02-E4} \\
BO067 & Résolution de l’équation du second degré par disjonction des cas (Démonstration) & C02-U2, p.\,\pageref{C02-U2} ; C02-EX2 & C02-E6, p.\,\pageref{C02-E6} \\
BO068 & Taux de variation et sécantes (Contenu) & C04-U1, p.\,\pageref{C04-U1} ; C04-EX1 & C04-E1, p.\,\pageref{C04-E1} \\
BO069 & Nombre dérivé comme limite intuitive du taux (Contenu) & C04-U2, p.\,\pageref{C04-U2} ; C04-EX2 & C04-E2, p.\,\pageref{C04-E2} \\
BO070 & Tangente comme limite des sécantes et pente (Contenu) & C04-U3, p.\,\pageref{C04-U3} ; C04-EX3 & C04-E5, p.\,\pageref{C04-E5} \\
BO071 & Approximation affine tangente (Contenu) & C04-U3, p.\,\pageref{C04-U3} ; C04-EX3 & C04-E5, p.\,\pageref{C04-E5} \\
BO072 & Calculer un taux et une pente de sécante (Capacité) & C04-U1, p.\,\pageref{C04-U1} ; C04-EX1 & C04-E1, p.\,\pageref{C04-E1} \\
BO073 & Interpréter pente, vitesse instantanée et coût marginal (Capacité) & C04-U2, p.\,\pageref{C04-U2} ; C04-EX2 & C04-E3, p.\,\pageref{C04-E3} \\
BO074 & Déterminer graphiquement un nombre dérivé (Capacité) & C04-U3, p.\,\pageref{C04-U3} ; C04-EX3 & C04-E4, p.\,\pageref{C04-E4} \\
BO075 & Construire une tangente de pente donnée (Capacité) & C04-U3, p.\,\pageref{C04-U3} ; C04-EX3 & C04-E4, p.\,\pageref{C04-E4} \\
BO076 & Déterminer l’équation de tangente (Capacité) & C04-U3, p.\,\pageref{C04-U3} ; C04-EX3 & C04-E4, p.\,\pageref{C04-E4} \\
BO077 & Calculer une approximation de f(a+h) (Capacité) & C04-U3, p.\,\pageref{C04-U3} ; C04-EX3 & C04-E4, p.\,\pageref{C04-E4} \\
BO078 & Équation de la tangente (Démonstration) & C04-U3, p.\,\pageref{C04-U3} ; C04-EX3 & C04-E2, p.\,\pageref{C04-E2} \\
BO079 & Racine carrée non dérivable en 0 (Démonstration) & C04-U4, p.\,\pageref{C04-U4} ; C04-EX4 & C04-AP4, p.\,\pageref{C04-AP4} \\
BO080 & Fonction dérivable et fonction dérivée (Contenu) & C06-U1, p.\,\pageref{C06-U1} ; C06-EX1 & C06-E1, p.\,\pageref{C06-E1} \\
BO081 & Dérivées carré, cube, inverse et racine (Contenu) & C06-U1, p.\,\pageref{C06-U1} ; C06-EX1 & C06-E1, p.\,\pageref{C06-E1} \\
BO082 & Dérivée de puissance entière positive ou négative (Contenu) & C06-U1, p.\,\pageref{C06-U1} ; C06-EX1 & C06-E1, p.\,\pageref{C06-E1} \\
BO083 & Somme, produit, inverse et quotient de fonctions dérivables (Contenu) & C06-U2, p.\,\pageref{C06-U2} ; C06-EX2 & C06-E1, p.\,\pageref{C06-E1} \\
BO084 & Valeur absolue : dérivabilité en 0 (Contenu) & C04-U4, p.\,\pageref{C04-U4} ; C04-EX4 & C04-AP4, p.\,\pageref{C04-AP4} \\
BO085 & Calculer une fonction dérivée par les opérations (Capacité) & C06-U2, p.\,\pageref{C06-U2} ; C06-EX2 & C06-E1, p.\,\pageref{C06-E1} \\
BO086 & Dérivée de la fonction carré (Démonstration) & C06-U1, p.\,\pageref{C06-U1} ; C06-EX1 & C04-E8, p.\,\pageref{C04-E8} \\
BO087 & Dérivée de la fonction inverse (Démonstration) & C06-U1, p.\,\pageref{C06-U1} ; C06-EX1 & C04-E8, p.\,\pageref{C04-E8} \\
BO088 & Dérivée d’un produit (Démonstration) & C06-U2, p.\,\pageref{C06-U2} ; C06-EX2 & C06-E1, p.\,\pageref{C06-E1} \\
BO089 & Fonctions paires et impaires, traductions graphiques (Contenu) & C06-U4, p.\,\pageref{C06-U4} ; C06-EX4 & C06-AP4, p.\,\pageref{C06-AP4} \\
BO090 & Lien signe de la dérivée et variations sur un intervalle (Contenu) & C06-U3, p.\,\pageref{C06-U3} ; C06-EX3 & C06-E6, p.\,\pageref{C06-E6} \\
BO091 & Caractérisation des fonctions constantes (Contenu) & C06-U3, p.\,\pageref{C06-U3} ; C06-EX3 & C06-E6, p.\,\pageref{C06-E6} \\
BO092 & Nombre dérivé et tangente à un extremum (Contenu) & C06-U3, p.\,\pageref{C06-U3} ; C06-EX3 & C06-E6, p.\,\pageref{C06-E6} \\
BO093 & Étudier variations et extrema (Capacité) & C06-U3, p.\,\pageref{C06-U3} ; C06-EX3 & C06-E2, p.\,\pageref{C06-E2} \\
BO094 & Étudier le second degré avec la dérivation (Capacité) & C06-U3, p.\,\pageref{C06-U3} ; C06-EX3 & C06-E2, p.\,\pageref{C06-E2} \\
BO095 & Résoudre une optimisation (Capacité) & C06-U5, p.\,\pageref{C06-U5} ; C06-EX5 & C06-E5, p.\,\pageref{C06-E5} \\
BO096 & Établir une inégalité (Capacité) & C06-U4, p.\,\pageref{C06-U4} ; C06-EX4 & C06-E3, p.\,\pageref{C06-E3} \\
BO097 & Étudier la position relative de deux courbes (Capacité) & C06-U4, p.\,\pageref{C06-U4} ; C06-EX4 & C06-E3, p.\,\pageref{C06-E3} \\
BO098 & Définition par f dérivée égale à f et f(0)=1 (Contenu) & C09-U1, p.\,\pageref{C09-U1} ; C09-EX1 & C09-AP1, p.\,\pageref{C09-AP1} \\
BO099 & Existence et unicité admises (Contenu) & C09-U1, p.\,\pageref{C09-U1} ; C09-EX1 & C09-AP1, p.\,\pageref{C09-AP1} \\
BO100 & Nombre e et notation puissance (Contenu) & C09-U1, p.\,\pageref{C09-U1} ; C09-EX1 & C09-AP1, p.\,\pageref{C09-AP1} \\
BO101 & Addition des arguments et produit (Contenu) & C09-U2, p.\,\pageref{C09-U2} ; C09-EX2 & C09-E1, p.\,\pageref{C09-E1} \\
BO102 & Inverse, quotient et puissances (Contenu) & C09-U2, p.\,\pageref{C09-U2} ; C09-EX2 & C09-E1, p.\,\pageref{C09-E1} \\
BO103 & Signe, sens de variation et courbe (Contenu) & C09-U2, p.\,\pageref{C09-U2} ; C09-EX2 & C09-E1, p.\,\pageref{C09-E1} \\
BO104 & Lien avec les suites géométriques (Contenu) & C09-U4, p.\,\pageref{C09-U4} ; C09-EX4 & C09-E8, p.\,\pageref{C09-E8} \\
BO105 & Transformer une expression exponentielle (Capacité) & C09-U2, p.\,\pageref{C09-U2} ; C09-EX2 & C09-E1, p.\,\pageref{C09-E1} \\
BO106 & Dériver exp(at) (Capacité) & C09-U3, p.\,\pageref{C09-U3} ; C09-EX3 & C09-E3, p.\,\pageref{C09-E3} \\
BO107 & Représenter exp(kt) et exp(-kt) pour k positif (Capacité) & C09-U4, p.\,\pageref{C09-U4} ; C09-EX4 & C09-E4, p.\,\pageref{C09-E4} \\
BO108 & Modéliser croissance ou décroissance exponentielle (Capacité) & C09-U4, p.\,\pageref{C09-U4} ; C09-EX4 & C09-E4, p.\,\pageref{C09-E4} \\
BO109 & Cercle trigonométrique, arc et radian (Contenu) & C10-U1, p.\,\pageref{C10-U1} ; C10-EX1 & C10-E1, p.\,\pageref{C10-E1} \\
BO110 & Enroulement de la droite et image d’un réel (Contenu) & C10-U2, p.\,\pageref{C10-U2} ; C10-EX2 & C10-E2, p.\,\pageref{C10-E2} \\
BO111 & Sinus, cosinus, triangle rectangle et valeurs remarquables (Contenu) & C10-U3, p.\,\pageref{C10-U3} ; C10-EX3 & C10-E3, p.\,\pageref{C10-E3} \\
BO112 & Placer un point sur le cercle (Capacité) & C10-U2, p.\,\pageref{C10-U2} ; C10-EX2 & C10-E2, p.\,\pageref{C10-E2} \\
BO113 & Lire les angles associés et leurs sinus et cosinus (Capacité) & C10-U4, p.\,\pageref{C10-U4} ; C10-EX4 & C10-E3, p.\,\pageref{C10-E3} \\
BO114 & Cosinus et sinus de $\dfrac{\pi}{4}$ (Démonstration) & C10-U3, p.\,\pageref{C10-U3} ; C10-EX3 & C10-RECH, p.\,\pageref{C10-RECH} \\
BO115 & Cosinus et sinus de $\dfrac{\pi}{3}$ (Démonstration) & C10-U3, p.\,\pageref{C10-U3} ; C10-EX3 & C10-RECH, p.\,\pageref{C10-RECH} \\
BO116 & Produit scalaire par projection et cosinus (Contenu) & C03-U1, p.\,\pageref{C03-U1} ; C03-EX1 & C03-E4, p.\,\pageref{C03-E4} \\
BO117 & Caractérisation de l’orthogonalité (Contenu) & C03-U1, p.\,\pageref{C03-U1} ; C03-EX1 & C03-E4, p.\,\pageref{C03-E4} \\
BO118 & Expression en base orthonormée et norme (Contenu) & C03-U2, p.\,\pageref{C03-U2} ; C03-EX2 & C03-E1, p.\,\pageref{C03-E1} \\
BO119 & Coordonnées par produits avec les vecteurs de base (Contenu) & C03-U2, p.\,\pageref{C03-U2} ; C03-EX2 & C03-E1, p.\,\pageref{C03-E1} \\
BO120 & Bilinéarité et symétrie (Contenu) & C03-U3, p.\,\pageref{C03-U3} ; C03-EX3 & C03-E3, p.\,\pageref{C03-E3} \\
BO121 & Carrés des normes de somme et différence (Contenu) & C03-U3, p.\,\pageref{C03-U3} ; C03-EX3 & C03-E3, p.\,\pageref{C03-E3} \\
BO122 & Formule d’Al-Kashi (Contenu) & C03-U4, p.\,\pageref{C03-U4} ; C03-EX4 & C03-E8, p.\,\pageref{C03-E8} \\
BO123 & Transformation de MA scalaire MB (Contenu) & C03-U5, p.\,\pageref{C03-U5} ; C03-EX5 & C03-E5, p.\,\pageref{C03-E5} \\
BO124 & Calculer angle et longueur (Capacité) & C03-U4, p.\,\pageref{C03-U4} ; C03-EX4 & C03-E3, p.\,\pageref{C03-E3} \\
BO125 & Choisir projection, coordonnées, normes et angle (Capacité) & C03-U4, p.\,\pageref{C03-U4} ; C03-EX4 & C03-E3, p.\,\pageref{C03-E3} \\
BO126 & Résoudre un problème géométrique (Capacité) & C03-U4, p.\,\pageref{C03-U4} ; C03-EX4 & C03-E3, p.\,\pageref{C03-E3} \\
BO127 & Démontrer une orthogonalité (Capacité) & C03-U2, p.\,\pageref{C03-U2} ; C03-EX2 & C03-E2, p.\,\pageref{C03-E2} \\
BO128 & Formule d’Al-Kashi par produit scalaire (Démonstration) & C03-U4, p.\,\pageref{C03-U4} ; C03-EX4 & C03-E8, p.\,\pageref{C03-E8} \\
BO129 & Lieu MA scalaire MB égal à zéro (Démonstration) & C03-U5, p.\,\pageref{C03-U5} ; C03-EX5 & C03-RECH, p.\,\pageref{C03-RECH} \\
BO130 & Vecteur normal et équation cartésienne (Contenu) & C07-U1, p.\,\pageref{C07-U1} ; C07-EX1 & C07-E1, p.\,\pageref{C07-E1} \\
BO131 & Projeté orthogonal sur une droite (Contenu) & C07-U2, p.\,\pageref{C07-U2} ; C07-EX2 & C07-E2, p.\,\pageref{C07-E2} \\
BO132 & Équation de cercle (Contenu) & C07-U3, p.\,\pageref{C07-U3} ; C07-EX3 & C07-E3, p.\,\pageref{C07-E3} \\
BO133 & Droite par un point et une normale (Capacité) & C07-U1, p.\,\pageref{C07-U1} ; C07-EX1 & C07-E1, p.\,\pageref{C07-E1} \\
BO134 & Coordonnées du projeté orthogonal (Capacité) & C07-U2, p.\,\pageref{C07-U2} ; C07-EX2 & C07-E2, p.\,\pageref{C07-E2} \\
BO135 & Équation à partir du centre et du rayon (Capacité) & C07-U3, p.\,\pageref{C07-U3} ; C07-EX3 & C07-E3, p.\,\pageref{C07-E3} \\
BO136 & Reconnaître cercle, centre et rayon (Capacité) & C07-U3, p.\,\pageref{C07-U3} ; C07-EX3 & C07-E3, p.\,\pageref{C07-E3} \\
BO137 & Utiliser un repère pour une configuration (Capacité) & C07-U4, p.\,\pageref{C07-U4} ; C07-EX4 & C07-E4, p.\,\pageref{C07-E4} \\
BO138 & Indépendance de deux événements (Contenu) & C05-U3, p.\,\pageref{C05-U3} ; C05-EX3 & C05-E3, p.\,\pageref{C05-E3} \\
BO139 & Partition de l’univers et probabilités totales (Contenu) & C05-U2, p.\,\pageref{C05-U2} ; C05-EX2 & C05-E2, p.\,\pageref{C05-E2} \\
BO140 & Succession de deux épreuves indépendantes (Contenu) & C05-U4, p.\,\pageref{C05-U4} ; C05-EX4 & C05-E5, p.\,\pageref{C05-E5} \\
BO141 & Répétition de Bernoulli identiques indépendantes pour n au plus 4 (Contenu) & C05-U4, p.\,\pageref{C05-U4} ; C05-EX4 & C05-E5, p.\,\pageref{C05-E5} \\
BO142 & Calculer par probabilités totales (Capacité) & C05-U2, p.\,\pageref{C05-U2} ; C05-EX2 & C05-E2, p.\,\pageref{C05-E2} \\
BO143 & Utiliser ou justifier l’indépendance (Capacité) & C05-U3, p.\,\pageref{C05-U3} ; C05-EX3 & C05-E3, p.\,\pageref{C05-E3} \\
BO144 & Arbre ou tableau de deux épreuves indépendantes (Capacité) & C05-U4, p.\,\pageref{C05-U4} ; C05-EX4 & C05-E5, p.\,\pageref{C05-E5} \\
BO145 & Arbre et calcul pour n au plus 4 Bernoulli (Capacité) & C05-U4, p.\,\pageref{C05-U4} ; C05-EX4 & C05-E5, p.\,\pageref{C05-E5} \\
BO146 & Variable comme fonction sur un univers fini (Contenu) & C08-U1, p.\,\pageref{C08-U1} ; C08-EX1 & C08-E1, p.\,\pageref{C08-E1} \\
BO147 & Loi de probabilité (Contenu) & C08-U1, p.\,\pageref{C08-U1} ; C08-EX1 & C08-E1, p.\,\pageref{C08-E1} \\
BO148 & Espérance (Contenu) & C08-U2, p.\,\pageref{C08-U2} ; C08-EX2 & C08-E5, p.\,\pageref{C08-E5} \\
BO149 & Linéarité de l’espérance (Contenu) & C08-U2, p.\,\pageref{C08-U2} ; C08-EX2 & C08-E5, p.\,\pageref{C08-E5} \\
BO150 & Variance et écart type (Contenu) & C08-U3, p.\,\pageref{C08-U3} ; C08-EX3 & C08-E2, p.\,\pageref{C08-E2} \\
BO151 & Formule de König–Huygens (Contenu) & C08-U3, p.\,\pageref{C08-U3} ; C08-EX3 & C08-E2, p.\,\pageref{C08-E2} \\
BO152 & Interpréter les événements X=a et X inférieur ou égal à a (Capacité) & C08-U1, p.\,\pageref{C08-U1} ; C08-EX1 & C08-E1, p.\,\pageref{C08-E1} \\
BO153 & Passer entre langage naturel et symboles (Capacité) & C08-U1, p.\,\pageref{C08-U1} ; C08-EX1 & C08-E1, p.\,\pageref{C08-E1} \\
BO154 & Modéliser par une variable (Capacité) & C08-U1, p.\,\pageref{C08-U1} ; C08-EX1 & C08-E1, p.\,\pageref{C08-E1} \\
BO155 & Déterminer une loi (Capacité) & C08-U1, p.\,\pageref{C08-U1} ; C08-EX1 & C08-E1, p.\,\pageref{C08-E1} \\
BO156 & Calculer espérance, variance, écart type (Capacité) & C08-U3, p.\,\pageref{C08-U3} ; C08-EX3 & C08-E2, p.\,\pageref{C08-E2} \\
BO157 & Utiliser l’espérance pour un jeu équitable (Capacité) & C08-U2, p.\,\pageref{C08-U2} ; C08-EX2 & C08-E3, p.\,\pageref{C08-E3} \\
BO158 & Simuler une variable (Capacité) & C08-U5, p.\,\pageref{C08-U5} ; C08-EX5 & TP, p.\,\pageref{TP} \\
BO159 & Écrire la moyenne d’un échantillon (Capacité) & C08-U5, p.\,\pageref{C08-U5} ; C08-EX5 & TP, p.\,\pageref{TP} \\
BO160 & Comparer moyenne simulée et espérance (Capacité) & C08-U5, p.\,\pageref{C08-U5} ; C08-EX5 & TP, p.\,\pageref{TP} \\
BO161 & Simuler N échantillons de taille n et fréquence dans 2 sigma sur racine n (Capacité) & C08-U5, p.\,\pageref{C08-U5} ; C08-EX5 & TP, p.\,\pageref{TP} \\
BO162 & Termes, sommes et seuils (Exemple algorithmique) & C01-U5, p.\,\pageref{C01-U5} ; C01-EX5 & TP, p.\,\pageref{TP} \\
BO163 & Factorielle, Fibonacci et Syracuse (Exemple algorithmique) & PYTHON, p.\,\pageref{PYTHON} ; PYTHON & TP, p.\,\pageref{TP} \\
BO164 & Liste des pentes des sécantes (Exemple algorithmique) & C04-U1, p.\,\pageref{C04-U1} ; C04-EX1 & TP, p.\,\pageref{TP} \\
BO167 & Approximation de pi par Archimède (Exemple algorithmique) & C10-U5, p.\,\pageref{C10-U5} ; C10-EX5 & TP, p.\,\pageref{TP} \\
BO168 & Monte-Carlo parabole et disque (Exemple algorithmique) & C05-U5, p.\,\pageref{C05-U5} ; C05-EX5 & TP, p.\,\pageref{TP} \\
BO169 & Espérance, variance et écart type en Python (Exemple algorithmique) & C08-U3, p.\,\pageref{C08-U3} ; C08-EX3 & TP, p.\,\pageref{TP} \\
BO170 & Fréquences des lettres (Exemple algorithmique) & C08-E7, p.\,\pageref{C08-E7} ; C08-E7 & TP, p.\,\pageref{TP} \\
\bottomrule\end{longtable}\normalsize